\clearpage
\thispagestyle{plain}
\phantomsection
\addcontentsline{toc}{chapter}{ABSTRACT}

{\centering\fontsize{12}{14}\selectfont\bfseries ABSTRACT\par}

\vspace{0.2cm}

% KONFIRMASI PENULIS: judul berbahasa Inggris adalah terjemahan \JudulPA, sesuaikan dengan judul resmi bila ada
{\noindent\raggedright
\textbf{IMPLEMENTATION OF A CONTINUOUS HEART RATE MONITORING SYSTEM BASED ON IOT EDGE COMPUTING FOR EARLY ARRHYTHMIA DETECTION}\par}

\vspace{0.2cm}

\noindent
\textbf{Name\hspace{1.55cm}: \NamaMahasiswa}\\
\textbf{NRP\hspace{1.75cm}: \NRP}\\
\textbf{Advisor\hspace{1.2cm}: \PembimbingSatu}\\
\textbf{\phantom{Advisor\hspace{1.2cm}: }\PembimbingDua}

\vspace{0.2cm}

Cardiovascular disease is the leading cause of death worldwide, and arrhythmia often occurs without symptoms, so continuous electrocardiogram (ECG) monitoring is needed for early detection. Existing monitoring systems are generally cloud-centric, with high latency and a dependence on internet connectivity. Arrhythmia detection models are also rarely tested under microcontroller constraints, and edge-to-cloud data pipelines are seldom tolerant of network disruption. This study builds a single-lead heart monitoring system based on edge computing using a TinyML approach. ECG preprocessing and arrhythmia classification run on-device on an ESP32-S3, and the results are sent through an MQTT pipeline designed to tolerate connection loss. A Butterworth filter and Pan-Tompkins R-peak detection run on the microcontroller, and a hybrid CNN with a morphology branch and a rhythm branch is quantized to INT8. Beats that have not been delivered are held in a local buffer in PSRAM. The float32 model on the MIT-BIH DS2 set, which contains 22 patients unseen during training, achieves a recall of 0.6996, an F1 of 0.6594, and an AUC of 0.9334. INT8 quantization yields a 22.94~KB model with a recall difference of $+$0.0108, and on-device inference takes 26.6~ms per beat. The C preprocessing is shown to be equivalent to the Python reference at 3.0\% of the time budget with no lost samples. The official body-signal recording test yields \PlaceholderData{number of recordings meeting the quality criteria out of the number of recording sessions}. Using a replayed MIT-BIH signal, the MQTT pipeline reaches a packet delivery success rate of 100.0\% and 99.2\% in two runs with median end-to-end latencies of 939 and 1,108~ms. Freezing the broker for 24 seconds loses no beat. The pipeline objective is achieved, and the preprocessing objective is achieved for the computation path while body-signal acquisition awaits the official recording test. The INT8 model objective is achieved except for the power consumption target, which has not been measured.
% sumber: terjemahan bab/abstrak.tex, angka identik dengan abstrak berbahasa Indonesia dan Bab 4 (desimal memakai titik sesuai kaidah bahasa Inggris)
% PLACEHOLDER-DATA abstract-akuisisi: terjemahan hasil pengujian rekaman sinyal tubuh resmi dari abstrak.tex (PLACEHOLDER-DATA abstrak-akuisisi)

\vspace{0.25cm}
\noindent
\textbf{Keywords:} ECG, IoT, Edge Computing, TinyML, ESP32-S3, Arrhythmia, MQTT

\clearpage
